\section{A Same-Object Audit Protocol}
\label{sec:protocol}

\begin{definition}[Source lock]
A candidate is a tuple
\[
  \mathcal C=(X,T,\tau,\varphi,\mathcal B,D),
\]
where \(X\) is the phase space, \(T\) the dynamics, \(\tau\) the clock or roof,
\(\varphi\) the potential and cocycle data, \(\mathcal B\) the function space,
and \(D\) the declared determinant convention.  The source lock also records
parameters, their provenance, numerical cutoffs, precision, allowed data, and
forbidden data.
\end{definition}

Changing any component that carries a failed obligation creates a new
candidate.  For example, adding reset edges to an acyclic level shift is not a
numerical regularization; it changes \(T\) and its periodic orbits.  Likewise,
assigning \(\log p\) to an orbit after looking up \(p\) changes the arithmetic
source and clock.

\paragraph{Five gates.}
The route decision is sequential: eligibility at a later gate requires the
same row to retain every earlier obligation.  The matrix nevertheless records
\emph{local structural support} at each gate.  Thus an S in A2 means that the
object has a valid determinant for its own declared ledger; it does not mean
that an object which failed A0 or A1 has passed the Route-A target.  Route
eligibility is the row-wise conjunction, never a column-wise assembly.
\begin{description}[leftmargin=3.1em,style=nextline]
  \item[A0: arithmetic origin.]
  Does the frozen rule generate the relevant rational-prime information and
  logarithmic scale without target-zero fitting or a prime-indexed lookup
  table?  A classical arithmetic grammar may still receive a weak or
  modeling-choice label when the primes are direct inputs rather than
  outputs.

  \item[A1: primitive/repetition ledger.]
  Are primitive objects, orientations, powers, multiplicities, and
  completeness intrinsic to the same dynamics?  A partition of finite states
  or a list of return atoms is insufficient unless it gives the declared
  periodic-orbit expansion.

  \item[A2: determinant.]
  Is a dynamical or Fredholm determinant defined on a named function space,
  with a trace expansion compatible with the A1 ledger?  A Dirichlet series
  or an inverse-designed analytic germ does not pass merely because it has an
  interesting quotient.

  \item[A3: global structure.]
  Are analytic continuation, functional equation, completed factors, total
  divisor growth, extra zeros and poles, and any proposed compression
  established for the same determinant?  A finite truncation cannot certify
  this gate.

  \item[A4: same-clock lift.]
  Is a Hilbert space, domain, and unitary, scattering, or self-adjoint
  construction defined while preserving the same clock, phases, and trace
  identity?  Formal analogy earns at most a partial label.
\end{description}

\paragraph{Evidence labels.}
We distinguish \evidence{proved}, \evidence{numerically certified},
\evidence{numerical observation}, \evidence{modeling choice},
\evidence{not testable}, and \evidence{open}.  Exact symbolic algebra and a
mathematical proof can receive the first label.  Exhaustive enumeration at a
frozen cutoff can certify the finite ledger it enumerates, but not an
asymptotic statement.  Floating-point orbit sums and cutoff trends remain
observations even when a precision audit is excellent.

For legibility, the figures compress the exact evaluator enums into S, P, and
F.  The compression is explicit and fail-closed: only named supported verdicts
map to S, only named weak/partial/formal verdicts map to P, and only named
failed/not-testable verdicts map to F.  An unrecognized enum stops figure
generation rather than receiving support by default.
\Cref{tab:candidate-matrix} also reports the first blocking gate and evidence
type for each row.

\paragraph{Adversarial controls and stop rules.}
Controls are chosen to test mechanism rather than visual similarity:
neighboring alphabets and roofs, finite-modulus approximants, on-circle and
off-circle polynomials, matched deletion rules, arithmetic and nonarithmetic
sign fields, and points on both sides of a proved convergence boundary.  If
the same inverse-design mechanism reconstructs both desired and adversarial
root geometries, the candidate stops as \evidence{proves too much}.  If an
early gate fails exactly, later attractive formulas are recorded but do not
unlock the route.

\paragraph{Data separation.}
No experiment loads a Riemann-zero table.  The Riemann--von Mangoldt law is
used only as a theorem-level growth benchmark \citep{dlmf2510}.  Rational
primes appear only when the candidate definition explicitly contains them or
when the candidate recursion generates them.  Every random seed is reported;
there is no best-seed selection.

\paragraph{Route lock.}
The subsequent operator evaluation, called \RouteB, is allowed only after the
same source-locked object reaches A4 readiness.  In the six records studied
here, all \RouteB{} flags are false.  This is a protocol outcome, not a
judgment about the existence of a Hilbert--Pólya program outside the frozen
objects.
